\chapter{Conclusions}
\label{ch:conclusions}

\section{Summary of Contributions}
\label{sec:summary_contributions}

This thesis set out to study whether recurrent spiking neural networks can
learn effectively in fully online reinforcement learning settings using
learning rules that rely only on local computations and a globally
broadcast reward signal, and made two main contributions toward answering
this question.

The first contribution is the implementation of a continuously running reinforcement-learning agent
based on an excitatory-inhibitory recurrent spiking neural network, trained online with the e-prop
framework \citep{eprop2020}, without episodic resets, on a GPU-accelerated, event-based simulation
substrate \citep{Yavuz2016GeNN}. This online, single-agent, no-reset formulation is a strict setting
for reward-based e-prop, and we did not find an existing implementation run under the same
conditions to compare against directly. As part of this implementation we also integrated
TD($\lambda$) into the e-prop update, decoupling the discount factor $\gamma$ from the credit-trace
decay $\lambda$, a formulation used throughout the experiments in Chapter~\ref{ch:results}.

The second contribution is \emph{distributed e-prop}, a new family of learning rules that eliminates
the structured feedback weights used by standard e-prop to route loss information back to each
neuron. Building on the theoretical connection between perturbation-based learning and policy
gradient methods, every neuron is instead treated as an independent stochastic actor that perturbs
its own membrane dynamics and correlates this local exploration with a single, identical, globally
broadcast temporal-difference error. The key novel ingredient that makes this work is the heuristic
introduction of the pseudo-derivative $\psi^t_j$ as a gating factor on the local signal, although the
perturbation derivation does not by itself call for it, borrowing $\psi^t_j$ from e-prop concentrates
each neuron's update on the moments when its membrane is near the firing threshold, where its
perturbations actually influence spiking, and suppresses it otherwise. This heuristic gating turns
out to be essential, removing it collapses learning and lets the membrane
voltage diverge (Chapter~\ref{ch:results}). Several variants of the rule were derived,
differing in where the perturbation is injected (at the neuron or at the synapse) and in whether the
local perturbation signal is additionally routed through the recurrent connectivity, and these
variants were evaluated against standard e-prop with adaptive feedback on
three reinforcement learning tasks of varying complexity.

\clearpage
\section{Main Findings}
\label{sec:main_findings}

The experiments in Chapter~\ref{ch:results} support several
findings that, taken together, characterize when and how each family of
learning rules succeeds or fails.

\paragraph{In the continuing, no-reset setting the reward rate is the natural measure of performance,
and the small discount factor the spiking rules need for a stable critic limits the \emph{policy}
they can learn with respect to it.} Using PPO only as a reference optimizer to probe the objective, the ceiling
experiments (Section~\ref{sec:ppo_ceiling}) show on Snake $5\times5$ that with a large
$\gamma = 0.99$ the episodic and continuing formulations lead to markedly different policies, one
favouring episode return and the other the reward rate, each poor on the other, whereas a small
$\gamma = 0.5$ reaches a good policy under neither, settling into a compromise that is suboptimal on
both metrics. Crucially, this is a limit on the achievable policy, not on the network's ability to
learn. The spiking rules
operate in exactly this small-$\gamma$ regime and reach roughly the operating point the reference
optimizer attains there, so on this task they sit near the ceiling \emph{for that regime}, what holds
them back is the regime itself, not their credit assignment. Decoupling $\lambda$ from $\gamma$
(Section~\ref{sec:practical-lambda}) is the natural way to try to lift this, keeping a small $\gamma$
for critic stability while extending the credit horizon through $\lambda$, but it is self-limiting:
as the same experiment and the analysis in Section~\ref{sec:ppo_ceiling} show, once the
critic predicts the relevant rewards well the higher-order terms a larger $\lambda$ would contribute
vanish, so the longer credit assignment is effectively switched off precisely where the critic is
accurate, and $\lambda$ can extend the policy's effective horizon only where the critic cannot
predict those rewards.

\paragraph{Symmetric e-prop fails as causal path length grows, and this
motivates the alternatives studied in this thesis.} On Snake $5\times5$,
symmetric e-prop was the strongest individual configuration tested, but
its performance collapsed entirely once per-head field layers were
introduced between the recurrent core and the output heads
(Section~\ref{sec:symmetric-failure}). Because its feedback weights are
tied directly to the forward output weights, any neuron not directly
connected to an output head receives no feedback signal whatsoever once
the causal path length exceeds one. This is not a gradual degradation but
an outright failure mode, and it directly motivates both random/adaptive
e-prop, whose feedback is decoupled from the forward weights, and
distributed e-prop, which removes structured feedback altogether.

\paragraph{Distributed e-prop is competitive with, and sometimes faster than, adaptive e-prop, but with
a lower asymptotic ceiling in some settings.} Across the Snake and Pac-Man comparisons
(Sections~\ref{sec:snake_comparison} and~\ref{sec:pacman_comparison}),
several distributed e-prop variants reached usable policies as quickly or more quickly than adaptive
e-prop within the fixed training budget used throughout this thesis. Over longer horizons, however,
the ordering can reverse. On the extended five-layer Snake network (Section~\ref{subsec:snake-depth-comparison}),
the distributed variants converged within the first few thousand episodes and then plateaued, while
adaptive e-prop, after an extended period of near-zero reward consistent with the time feedback
alignment takes to take effect, continued improving past the distributed variants' plateau for the
remainder of the roughly $100{,}000$-episode run. The Connect Four results
(Section~\ref{sec:connect_four}) are consistent with the same ordering, adaptive e-prop ends up
highest and still slowly climbing while the distributed variants stay flat, although on that task the
runs differ in length, no rule reaches a usable policy, and adaptive's continued gain coincides with
steadily worsening spiking dynamics. Neither rule
should therefore be read as uniformly faster or uniformly better; their relative standing depends on
both the training budget and the task at hand.

\paragraph{Neither learning rule is uniformly more robust to the choice of
$\lambda$, and the two tasks studied gave opposite answers.} On Snake
$5\times5$ (Section~\ref{sec:snake_lambda_comparison}), adaptive e-prop
tolerated nearly the full range of $\lambda$ tested, failing only at the
most extreme value, while the distributed variant showed a narrower and
less monotonic band of values for which it learned promptly. On the
Pac-Man randomized layouts problem
(Section~\ref{sec:pacman_comparison}), the
distributed variant tolerated a wide range of $\lambda$, while adaptive
e-prop showed an inconsistent and often delayed learning onset across the
same sweep. This reversal is itself one of the more informative results
in this thesis, since it demonstrates directly that robustness to the
credit-assignment horizon is not an intrinsic property of either
algorithm, but depends on the task, and possibly on other interacting
factors such as architecture depth or reward sparsity that differ between
the two environments.

\paragraph{Distributed e-prop exhibits at least two distinct failure modes
under a poorly chosen $\lambda$, not one.} On the Pac-Man sweep, large
values of $\lambda$ destabilized the network's spiking dynamics, but not
always in the same direction, some values drove the network toward
runaway, pathologically high firing, while at least one drove it toward
the opposite extreme of excessive silence. Both are consistent with the
underlying credit trace losing useful information at long horizons, but
the fact that the same learning rule can fail in either direction, under
broadly similar conditions, suggests that the relationship between
$\lambda$ and network-level stability in this family of rules is more
intricate than a simple variance argument would predict.

\paragraph{The perturbation-based structure of distributed e-prop enables a
form of local reward shaping that standard e-prop could not support in
the same way.} The voltage-regularization experiment
(Section~\ref{sec:regularization-exploration}) found that a local,
perturbation-gated voltage reward could be used to keep membrane
potentials in a well-regulated range for distributed e-prop, while the
analogous regularization term applied within standard e-prop caused
membrane voltages to diverge instead of stabilizing it. The likely explanation
is that the perturbation trace naturally carries little signal when the
membrane is far from threshold, gating the regularization update away
from precisely the regime in which it would otherwise be poorly directed,
a self-suppression property that the standard e-prop update lacks.

\paragraph{Distributed e-prop is markedly more robust to a violated
modeling assumption in the neuron model than adaptive e-prop in at least one
observed setting.} When
membrane integration was disabled during the refractory period, a change
that invalidates an assumption in the analytical derivation of standard
e-prop's eligibility trace, adaptive e-prop failed to learn at all on the
Pac-Man randomized layouts problem, while the distributed variants
continued to learn successfully and maintained well-regulated spiking
dynamics (Section~\ref{subsec:no_refractory_integration}). This is consistent
with distributed e-prop's credit assignment being based on perturbation
correlations, which makes it comparatively insensitive to violations of the
specific assumptions that decomposition relies on.

\paragraph{The heuristic pseudo-derivative gating is essential to distributed e-prop, not an optional
refinement.} Although the perturbation-based derivation does not itself require it, gating each
neuron's local update by the pseudo-derivative $\psi^t_j$, so that credit accrues chiefly when the
membrane is near the firing threshold, proves necessary for the rule to learn at all. The ablation
that removes this factor (\emph{weight dist no heuristic}, Section~\ref{sec:snake_comparison}) does
not merely learn more slowly ,its performance collapses well below every other recurrent variant and
its membrane voltage diverges, growing without bound. This identifies the heuristic gating as a load-
bearing component of the rule, and suggests
that concentrating credit on the threshold-crossing regime, where a neuron's perturbations actually
affect its spiking, is what keeps the perturbation-based update both well-directed and dynamically
stable.

\paragraph{Distributed e-prop learns representations from which the value head regresses faster and more
cleanly than adaptive e-prop does, even though it carries no information about how it affects the
value function, which is one of the most surprising results of this work.} Across the qualitative
value-trace comparisons (Section~\ref{sec:value-traces}, Figure~\ref{fig:value_function_comparison}),
the distributed variants develop sharper, more anticipatory value estimates earlier in training than
adaptive e-prop, indicating that their recurrent representations are more readily linearly decodable
into an accurate value signal. This is counterintuitive given the two rules' information content:
adaptive e-prop routes an approximation of the loss gradient back through its feedback weights, so it
carries an approximate value gradient, an explicit signal of how each neuron should change to reduce
the value-prediction error, whereas distributed e-prop has no analogue at all, assigning credit only
from the correlation between local perturbations and a single broadcast scalar. On the value-function
regression task specifically, the gradient-aligned rule would therefore be expected to hold the
advantage, and yet the perturbation-based rule produces the more value-friendly representations. The
result suggests that shaping a representation that supports value estimation does not require an
explicit value gradient, and that the perturbation-based dynamics confer some advantage on
representation learning that the structured-feedback rule does not, an effect worth isolating in
future work.

\paragraph{The learning rules studied here are genuinely different objects, not minor variants of one
another, and distributed e-prop in particular performs surprisingly well given how much it gives up,
which makes it a promising target for further study.} Throughout the experiments the rules differed
not only in final performance but in how they learned, in their spiking dynamics, their sensitivity
to depth and to $\lambda$, their failure modes, and their robustness to violated modeling
assumptions, so they are better understood as distinct credit-assignment mechanisms that each
warrant study in their own right than as interchangeable choices. What stands out most is that
distributed e-prop remains competitive at all under the strictest constraints considered in this
thesis, it discards the structured feedback weights of standard e-prop entirely, assigns credit from
nothing more than the correlation between each neuron's own local perturbation and a single,
globally broadcast scalar TD error, and does so fully online and in constant memory. That a rule this
unconstrained can match or, in some settings, outpace adaptive e-prop, tolerate a violated neuron
model that defeats it, and support local reward shaping it cannot, is the most surprising outcome of
this work. It strongly suggests that distributed e-prop is worth studying further, both to map where
its advantages hold and to find the tasks and regimes that expose its limits.

\paragraph{High firing rates are a consistent operating regime, not a
pathology of any single rule.} Across all experiments, all learning rules, and
all three tasks, firing rates converge to a band of roughly 130--170~Hz,
far above both the biological range ($\approx$5--20~Hz) and the target rate
$F^* = 10$~Hz used in the firing-rate regularization. This holds regardless of
credit-assignment strategy, network depth, or task. The firing-rate
regularization term imported from standard e-prop proves insufficient to keep
rates near the target. Our hypothesis is that the some of the tasks studied here 
(Snake, Pac-Man) modeled as they are in this work, 
are highly reactive, requiring the network to maintain near-constant 
firing to produce value and policy estimates at every
step, which pushes the dynamics toward dense activity. Additionally, when the
learning signal is weak or absent, the dynamics tend to drift toward
saturation, as the mechanisms that would otherwise
reinforce sparse, informative firing patterns lack the signal to do so.
Whether stronger rate-control mechanisms or tasks with different temporal
structure would allow the network to settle into a sparser regime remains an
open question.

\paragraph{A single, globally broadcast scalar signal is enough to train recurrent spiking networks
effectively online, though with qualifications.} The question of whether credit assignment from
nothing more than local computation and one identical, globally broadcast reward-related signal can
drive effective learning, recurring throughout this thesis, is answered here in the affirmative for
the tasks studied: the distributed e-prop variants learn non-trivial behavior on Snake and Pac-Man,
including the anticipatory, partially observable navigation the latter requires, and do so fully
online, in constant memory, and within training budgets comparable to, and at times shorter than,
those of the gradient-aligned adaptive rule. The qualifications are equally part of the finding. The
extended-training comparison shows the distributed variants plateauing below adaptive e-prop's
eventual ceiling (Section~\ref{subsec:snake-depth-comparison}), the usable range of $\lambda$ is
task-dependent and can destabilize the dynamics when poorly chosen, and the hardest, maximally
sparse task, Connect Four, was not solved by any rule (Section~\ref{sec:connect_four}). The
broadcast signal is therefore sufficient to learn fairly complex policies efficiently, but not
without limits on asymptotic performance and on the difficulty of task it can currently reach.

\section{Limitations}
\label{sec:limitations}

Several limitations should be kept in mind when interpreting the results
of this thesis.

\paragraph{Most comparisons rely on a small number of runs, and several of
the depth and $\lambda$ sweeps reported in
Chapter~\ref{ch:results} are single runs per configuration.} Where
this is the case, it has been noted explicitly alongside the relevant
results, but it means that some of the more striking individual
observations, such as the specific failure of $\lambda = 0.4$ for the
\texttt{combined} variant on Snake, or the specific failure mode taken by
$\lambda = 0.9$ for \texttt{weight\_dist\_prop} on Pac-Man, cannot be
distinguished from idiosyncrasies of that particular run without repeated
seeds. The five-run comparisons presented for the primary Snake and
Pac-Man evaluations are more robust in this respect, but even there, the
hyperparameters were not tuned exhaustively for any single learning rule,
and a configuration that favors one rule may be poorly matched to
another, as noted repeatedly throughout Chapter~\ref{ch:results}.

\paragraph{No direct comparison was made against Truncated Backpropagation Through
Time as upper-bound baseline for task
performance.} The motivation for this thesis is squarely the online,
constant-memory setting in which this method is not applicable
in its unmodified form, but this also means that the absolute
performance ceiling achievable on these tasks is not known for the arquitectures
tested in this thesis. If the requirements were relaxed, 
we could test how these architectures would work using TBPTT and test it against 
the presented rules.

\paragraph{Several proposed extensions of distributed e-prop did not work
reliably enough to be included in the main comparisons.} The per-synapse
exploration-variance update (Section~\ref{subsec:exploration_variance}) was
found to add complexity without a clear benefit over the simpler
neuron-level update, and was therefore kept fixed in the main
experiments. The propagation of local regularization rewards through the
structured-routing machinery (Section~\ref{subsec:regularization}) showed
promising preliminary results but was not investigated systematically.

\section{Future Work}
\label{sec:future_work}

The limitations and open questions identified throughout this thesis
suggest several directions for future work, ranging from immediate
methodological extensions to more open-ended questions about how local,
perturbation-based learning rules might be developed further.

\paragraph{Systematic characterization of the discount $\gamma$.}
Throughout this work the discount was fixed to $\gamma = 0.5$ on the basis of
initial empirical exploration. This was
sufficient to obtain stable learning, but it leaves the role of $\gamma$
largely uncharted in our setting. The decoupling introduced in
Section~\ref{sec:decoupled-lambda-rule} now makes such a study well-posed, since
$\gamma$ is confined to the critic and no longer constrains the actor's credit
horizon, its effect is cleanly isolated to value estimation, and a dedicated
$\gamma$-sweep can be interpreted without confounding it with the
action-credit-assignment range. A natural next step is therefore to map, for
each learning rule, how critic learnability degrades as $\gamma$ grows,
identifying the point at which the compounding of small per-step value errors
across bootstrapping destabilizes the target $\delta_t = R_{t+1} + \gamma V_{t+1} -
V_t$ and the value head ceases to converge. Because the learning rules differ
in how they route the global feedback signal, their critics may tolerate
different maximum discounts; characterizing this per-rule threshold would
clarify whether the aggressive discount required here is a fundamental
constraint of the continuing, no-reset setting with a noisy spike-based critic,
or an artifact specific to certain rules. The interaction with an optional
critic-specific $\lambda_V$ and with the absence of episodic resets would be
worth examining in the same sweep.

\paragraph{Combining structured feedback and perturbation-based credit
assignment more deliberately.} Distributed e-prop was originally
conceived as a way of correcting the bias introduced by the approximate
feedback weights of standard e-prop, by supplementing the
structured learning signal with a perturbation-based estimate. In
practice, the \texttt{combined} variant explored in
Chapter~\ref{ch:results}, which simply sums the synaptic
perturbation update with the adaptively routed feedback update, behaved
similarly to the perturbation-only variants and did not show a clearly
synergistic improvement over either rule in isolation. This is somewhat
surprising given the flexibility of the underlying framework, and it
seems unlikely that no combination of the two signals can outperform
each individually. More structured ways of combining them remain
unexplored, the two signals could be weighted dynamically, 
for instance based on the agreement between the two estimates, 
in a manner similar to the heuristic agreement 
term already introduced for the perturbation and
backrouted traces in Section~\ref{subsec:structured_routing}; the feedback
signal could be used primarily as a variance-reduction term early in
training, when the perturbation-based estimate is least reliable, and
phased out as the feedback weights converge; or the two updates could be
applied to disjoint subsets of weights. 
Identifying a combination strategy that reliably outperforms
both constituent rules would be a natural next step.

\paragraph{Ablating the role of learned output weights.} As discussed in
Section~\ref{subsec:neuron_as_actor}, the output weights receive an exact
gradient at every step, unlike the perturbation-based estimates that drive the
recurrent layer, and the quality of the readout they learn largely determines
how informative the broadcast error $`\delta_t`$ is for the hidden neurons.
This suggests that much of the early learning may be driven by the readout
adapting to the recurrent layer's initial dynamics. A dedicated ablation
study (freezing the output weights at random initialization, or delaying
their learning for a fixed number of steps) would help disentangle how much
of the learning is attributable to the recurrent rule versus to the
adaptation of the readout that supplies its credit signal.

\paragraph{Alternative gating functions for the pseudo-derivative modulation.}
The pseudo-derivative $`\psi_j^t`$ that gates the perturbation-based update
is a specific choice, not a theoretical necessity, as noted in
Section~\ref{subsec:lif_update_rule}, one could replace
it with alternative gating mechanisms, such as smoother functions of the
membrane potential, and still obtain valid variants of the update rule. The
ablation in Chapter~\ref{ch:results} showed that removing the gating
altogether collapses learning, but this does not establish that the
particular form of $`\psi_j^t`$ borrowed from e-prop is optimal. A
systematic study of alternative gating functions (for instance smooth
sigmoid-like thresholds with tunable width, or functions that depend on
the recent spiking history of the neuron) could reveal whether the current gating is merely
sufficient or close to best among a broader family, and whether the
threshold-crossing regime it concentrates credit on is the most informative
regime for perturbation-based updates in general.

\paragraph{Alternative routing weights for the distributed routing pathway.}
The routed variants of distributed e-prop reuse the recurrent weights
$`W^{\mathrm{rec}}`$ themselves to propagate each neuron's local signal through
the network (Section~\ref{sec:feedback_connectivity}), which makes the routing
symmetric by construction, tied directly to the forward connectivity. This is
only one of the choices available, and it mirrors symmetric e-prop. The same
decoupling explored for standard e-prop's feedback matrices could be applied
here as the distributed routing could instead use a fixed random projection or an
adaptively learned matrix, separate from $`W^{\mathrm{rec}}`$, yielding random
and adaptive analogues of the routed distributed variants. Studying these would
clarify whether the benefits of routing the local signal stem specifically from
its alignment with the forward recurrent weights, or whether any structured
projection of the local signals across the population suffices, paralleling the
question of how much symmetric e-prop's performance depends on weight alignment
versus merely having a structured feedback path at all.

\paragraph{Exploring a gated per-neuron exploration-variance update.} The
per-neuron exploration variance $`\sigma_j`$ was held fixed throughout the
main experiments, and the adaptive update proposed in
Section~\ref{subsec:exploration_variance} did not yield a clear benefit over
a fixed value. One possible reason is that the variance update, as
formulated, accumulates credit from perturbations regardless of whether the
neuron's membrane was near threshold at the time, even though it is
precisely in that near-threshold regime that the perturbation actually
influences spiking. The pseudo-derivative gating $`\psi_j^t`$ was shown to
be essential for the weight update for exactly this reason, but the
variance update has no analogous mechanism. A natural direction is to
apply a similar gating factor to the $`\sigma_j`$ update, so that the
variance is adapted predominantly from perturbation correlations at
moments when the neuron's perturbations were behaviorally relevant. This
could reduce the noise in the variance estimate itself and potentially make
the adaptive $`\sigma_j`$ update viable where the ungated version was not.

\paragraph{Richer and state-dependent exploration.} The exploration variance
$\sigma_j$ used throughout this thesis is either a fixed hyperparameter
or, at best, a per-neuron scalar updated by its own local policy gradient
(Section~\ref{subsec:exploration_variance}), with no dependence on the
network's current state. A natural extension is to make $\sigma_j$ a
function of context, for instance increasing exploration in states where
the value estimate is uncertain or rapidly changing, and decreasing it
once a stable policy has been found for a given region of the state
space. This connects directly to ideas from meta-learning, since a
state-dependent exploration policy is itself a learned controller that
shapes how the underlying weights are updated, an instance of the network
modulating its own learning; we return to
this connection below. A second, largely orthogonal direction concerns
the structure of the injected noise itself. The perturbation variants
studied in this thesis inject either fully independent noise per neuron
or per synapse, or a single shared sample per postsynaptic neuron scaled
individually per synapse (Section~\ref{subsec:synaptic_perturbation}).
An interesting intermediate possibility, not explored here, is to draw
the perturbation noise from a small shared latent space, with a low-
dimensional noise vector sampled once per timestep and projected onto
each neuron or synapse through a fixed or learned projection. This would
deliberately correlate the exploration noise across the population,
in a manner reminiscent of how evolution strategies can benefit from
structured perturbations, and could offer
a different point in the bias-variance tradeoff than the perturbation
variants compared in Chapter~\ref{ch:results}.

\paragraph{A broader family of local reward signals.} The voltage- and
firing-rate-regularization rewards introduced in
Section~\ref{subsec:regularization} demonstrated that distributed e-prop can
incorporate purely local evaluative signals, compatible with its
perturbation-based update, in a way that standard e-prop could not
support as effectively. Voltage and firing rate are only two examples of
quantities a neuron could evaluate about itself; other candidates might
include the temporal sparsity or burstiness of a neuron's own spike
train, the correlation of a neuron's activity with that of its immediate
neighbors, or measures of how informative a neuron's recent activity has
been for the network's eventual output. More broadly, this raises the
question of what other local objectives admit a similarly natural
perturbation-based update, and whether a network exposed to several such
local rewards simultaneously, alongside the global task reward, would
develop qualitatively different internal representations than one trained
on task reward alone.

\paragraph{Structural plasticity and online rewiring.} Throughout this thesis the
network's synaptic topology was fixed at initialization, the toroidal Gaussian
connectivity, the excitatory-inhibitory partitioning, and the per-head field
layers were all static for the duration of training. Only the synaptic weights
were adapted. Yet in biological circuits, learning is not confined to changing
existing connection strengths; structural plasticity, including synaptic
formation and elimination, rewires the graph itself on a slow timescale, and
this process is thought to play a role in stabilizing dynamics and expanding
the functional representations available to the agent. An analogous mechanism
could be valuable in the setting studied here, particularly given the finding
that all learning rules converge to a high-firing regime with poorly
regulated membrane voltages (Section~\ref{sec:main_findings}).

\paragraph{Learning signals produced by a self-trained mechanism.} Throughout this thesis, the learning
signal each neuron receives is either a structured but fixed
approximation to the gradient, routed through feedback weights that are
random, symmetric, or themselves adaptively learned
(Section~\ref{subsec:eprop-variants}), or a single scalar TD error, identical
for every neuron, in the case of distributed e-prop.\
\citep{OneShotEprop} explore a related but distinct idea: a separate
network is trained to produce the per-neuron learning signal $L^t_j$
itself, taking on a role loosely analogous to dopaminergic areas of the
brain that gate plasticity elsewhere.

One possibility is a population of additional, self-trained neurons that
act as a kind of student network, learning online, alongside the main
network, to predict or shape the contribution each neuron should make to
the eventual outcome. A second, less architecturally
involved possibility is to give each neuron a supplementary learned
signal, produced by an additional set of
weights. A third possibility, echoing
the shared-latent-space idea discussed above for exploration noise, is to
project a small, shared, learned latent vector onto each neuron at every
timestep, so that neurons receive a population-level signal that is
richer than a single scalar but still far more constrained than a fully
structured, per-neuron feedback projection.

A point worth stressing for all three options is that, for them to remain
consistent with the rest of this thesis, whatever additional mechanism
produces these signals must itself be trained online, using the same
continuous stream of experience and the same constant-memory,
forward-in-time updates as the rest of the network. The appeal of distributed
e-prop is precisely that every component of its update is local and
online; a supplementary signal-generating mechanism that itself depended
on offline meta-training, or on machinery incompatible with the
online setting motivating this thesis, would reintroduce exactly the
kind of constraint the rest of this work tries to avoid. Whether a
self-trained signal of this kind, learned under the same online
constraints as everything else, can meaningfully enrich the credit
assignment available to each neuron without compromising locality is an
open question this thesis leaves for future investigation.

\paragraph{A more systematic exploration of architecture size and depth.}
The depth comparison in Section~\ref{subsec:snake-depth-comparison} varied the
number of layers for two specific learning rules on a single task, and
found that the two rules respond differently to depth within a fixed
training budget. This leaves open how either rule behaves as other
architectural dimensions are varied, such as the width of each
excitatory-inhibitory population, the fan-in of the toroidal Gaussian
connectivity, or the presence and size of the per-head field layers
independently of depth. A more systematic sweep across these dimensions,
ideally with the repeated seeds discussed above, would help separate
effects that are specific to causal path length from effects that are
more generally about network scale or capacity.

\paragraph{Tasks demanding longer-term memory.} The three environments
studied in this thesis were chosen to probe sparse-reward credit
assignment, working memory under partial observability, and behavioral
adaptation, but none of them require the agent to retain information over
very long horizons relative to the timescales of the underlying neuron
and adaptation dynamics. Tasks that demand longer-term memory, for
instance environments where relevant information must be retained for
hundreds or thousands of timesteps before it becomes useful, would test
the limits of the adaptive threshold mechanism's slow timescale 
and of the action-credit-assignment
horizon controlled by $\lambda$ considerably more severely than the tasks
studied here. Such tasks would also raise practical questions about how
training itself should be adapted, since a longer required memory horizon
may demand a correspondingly longer $\lambda$, which Chapter~\ref{ch:results}
showed can itself destabilize the algorithms spiking dynamics;
resolving this tension may require some of the other directions
discussed above, such as a state-dependent or self-modulated $\lambda$.

\paragraph{Repeated-seed evaluation of the depth and $\lambda$ sweeps.} The
single-run nature of several of the comparisons in this thesis, noted in
Section~\ref{sec:limitations}, means that converting the more striking
qualitative observations, such as the divergent failure modes of
distributed e-prop at large $\lambda$ on Pac-Man, into statistically
grounded claims would require repeating these sweeps across multiple
seeds.

\paragraph{A systematic study of the relationship between $\lambda$ and
network-level stability in distributed e-prop.} The observation that
large $\lambda$ can drive the network toward either runaway firing or
excessive silence, depending on the specific value and possibly the
random seed, is currently unexplained. Understanding what determines
which failure mode a given configuration falls into would clarify
whether this is a fundamental property of the perturbation-based update
or an artifact of the specific architecture and task studied here.

\paragraph{Deployment on neuromorphic hardware.} The implementation
developed in this thesis was deliberately constrained to be expressible
entirely in terms of local, per-synapse updates and event-driven
computation (Section~\ref{sec:implementation}), with the explicit
goal of acting as a faithful prototype for spike-driven accelerators such
as Loihi \citep{Davies2018Loihi} or SpiNNaker \citep{Painkras2013SpiNNaker}.
Carrying out this deployment directly would be a natural next step, and would allow
the energy-efficiency motivations discussed in Chapter~\ref{ch:introduction} to
be evaluated empirically.
